% Atur variabel berikut sesuai namanya

% nama
\newcommand{\name}{Yoel Yhokhanan Sianipar}
\newcommand{\authorname}{Sianipar, Yoel Yhokhanan}
\newcommand{\nickname}{Yoel}
\newcommand{\advisor}{Prof. Dr. Muhammad Rivai, S.T., M.T.}
\newcommand{\coadvisor}{Dr. Suwito, S.T., M.T.}
\newcommand{\examinerone}{Fajar Budiman, S.T., M.Sc.}
\newcommand{\examinertwo}{Ir. Tasripan, M.T.}
\newcommand{\examinerthree}{}
\newcommand{\headofdepartment}{Ronny Mardiyanto, S.T., M.T., Ph.D.}

% identitas
\newcommand{\nrp}{5022221081}
\newcommand{\advisornip}{196904261994031003}
\newcommand{\coadvisornip}{198101052005011004}
\newcommand{\examineronenip}{198607072014041001}
\newcommand{\examinertwonip}{196204181990031004}
\newcommand{\examinerthreenip}{}
\newcommand{\headofdepartmentnip}{198101182003121003}

% judul
\newcommand{\tatitle}{SISTEM PENGENDALIAN TEKANAN POMPA AIR BERBASIS INVERTER DENGAN \emph{NEURAL NETWORK}}
\newcommand{\engtatitle}{\emph{INVERTER-SUPPLIED WATER PUMP PRESSURE CONTROL SYSTEM USING NEURAL NETWORK}}

% tempat
\newcommand{\place}{Surabaya}

% jurusan
\newcommand{\studyprogram}{Teknik Elektro}
\newcommand{\engstudyprogram}{Electrical Engineering}

% fakultas
\newcommand{\faculty}{Teknologi Elektro dan Informatika Cerdas}
\newcommand{\engfaculty}{Intelligent Electrical and Informatics Technology}

% singkatan fakultas
\newcommand{\facultyshort}{FTEIC}
\newcommand{\engfacultyshort}{FTEIC}

% departemen
\newcommand{\department}{Teknik Elektro}
\newcommand{\engdepartment}{Electrical Engineering}

% kode mata kuliah
\newcommand{\coursecode}{EE234799}
